% !TeX root = ../papers/nc-p3.tex
\section{有限量子ファイバーの三つの六次元正規形}
\label{sec:finite-quantum-planar-fibers}
\begin{theorem}\label{thm:finite-quantum-planar-fibers}
四辺が有限で量子閉極限を持ち、三次元接点がない全窓の添字保持類は、次の互いに素な三表で尽きる。
各表で異なる六座標は非同型であり、全点の射影安定群は自明である。
\begin{enumerate}
\item 共通面値 $-1$ では $(a,b,c,d,e,f)\in\mathbb A^6$ の一表である。
\item 原始四乗根の閉類では、二つの向きにそれぞれ
$\mathbb A^6\setminus\{b=c=d=e=f=0\}$ が一表ずつある。
除いた直線の窓も全次数で正則であり、そこでは点対応に曲線が残る。
\end{enumerate}
向きを選んで $\zeta=-1,i,-i$ と書くと、一移動又はその逆は六座標を順に
$\zeta^{2},\zeta^{3},\zeta^{3},1,1,\zeta$ 倍する。最小周期は非零座標に対応する倍率の位数の最小公倍数である。
前節の最大接点の分類と合わせて、二つの量子ファイバーにある有限グラフを全て扱っている。
\end{theorem}
\begin{proof}[余重みと印なし同型]
前定理で最大接点がないとき、余重みは $w=(0,3,5,9)$ で三係数 $A,B,C$ は非零である。
背景の二文字交換倍率は位置0の文字を他と可換として、
残る $12,13,23$ で $\zeta,\zeta^{-1},\zeta$ に取る。
体積の順列係数は符号に各反転対のこの倍率を掛ける（逆の規約は向きを反転する）。
多重次数の差は四添字の和が零の格子に取る。
三つの余重み1の方向を
\[
 (-1,0,2,-1),\quad(-1,2,-1,0),\quad(1,-1,-1,1)
\]
の順に取れば格子の整数基底である。よって三係数は共通対角群で指定した非零値に一意にできる。
正の残る自由方向は次表で、項目は多重次数、余重み、座標の順である。
\[
\begin{array}{c|r|c@{\qquad}c|r|c}
(-1,1,-1,1)&7&a&(-1,-1,2,0)&7&b\\
(0,-1,-1,2)&10&c&(-1,-1,1,1)&11&d\\
(-1,0,-1,2)&13&e&(-1,-1,0,2)&15&f
\end{array}
\]

具体的な復元を述べる。定理\ref{thm:quantum-universal-chambers}のSchur補の二切断を使う。
各重みで語を辞書順に並べ、線型写像の核は既約行階段形の自由変数を順に1とする基底に取る。
非斉次式の特解には、この核と標準Hermite内積で直交する唯一の解を取る。
ここで複素共役は定数 $\mathbb Q(i)$ のみで使い、パラメータには作用させない。
重み1の核の三座標は全て1とする。表の六方向では座標を加え、他の核方向は軌道接なので零とする。
$\zeta=i$ では正規形の $d,e$ に順に $504i/78125,-7911i/12500000$ を加えて使う。
反対の向きでは共役する。余重み19まで行えば全256係数が定まる。

\path{computations/verify_finite_quantum_planar_fibers.py} が各重みで二つの
Schur補の右辺が線型部分の像に入ることを六変数の多項式恒等式として検算する。
実装では主六行列の逆の次数 $j>0$ 成分を $I_j=-I_0\sum_{l=1}^j M_{l,PP} I_{j-l}$ で得る。
主行列式は重み零で一定なのでこの有限再帰で評価できる。核の負添字二個の成分は記載した九線、
他の非零核は全て根指標である。根指標の方向は二つの面値1を要する横断変形がなく
（定理\ref{thm:quantum-face-resonance}の外積計算）、全て軌道接となる。
したがって直交条件は実際のLuna横断面を取っている。
背景の共通トーラスと有限単項式安定群、巡回作用はこの補空間を保つ。
実際コンパクト対角部分と単位絶対値の置換持上げはユニタリで、対角部分は稠密である。
これで多項式再帰のパラメータは飽和スライス内の軌道の代表を与える。

上の三基底を正にする余指標の開錐は、最低を0とすると正確に
$e_2<2e_1$, $e_1+e_2<e_3<2e_2$ である（ここだけ $e_j$ は矢重みを示す）。
この錐の内部では常に $0<e_1<e_2<e_3$ である。表の六方向は三基底の非負整数和で、
再帰の他の成分もそれらの和のみから生じる。従って流入する余指標の開錐は
三非零一次方向だけで既に決まり、どのパラメータでも同じになる。
この錐から文字の順が復元されるため、同じ順を保つ置換は恒等である。
背景安定群の置換は $-1$ では全対称群、原始四乗根では交代群なので、
順の同値類はそれぞれ一つ、二つである。トーラスにも三単元を固定する非自明な元はない。
Lunaのファイバー同値から元の窓でも安定群と同型は記載通りである。
巡回は正規化前の各核線で根付き単位倍率を与え、指標基底の三倍率を共通対角群で戻す。
上記計算は全テンソルに対するこの巡回と六座標の指定倍率を直接比較する。
\end{proof}

\begin{proof}[有限性の全特殊値]
辺の二端を $[z:u:v:H],[z':u':v':H']$ とする。$H=0$ なら $r_{23},r_{13}$ から
$H'=0$ 又は $u=v=0$ となり、後者では最低点の相手も最低点である。逆も同じである。
$H=H'=0$ では $r_{03}$ が $Cvv'=0$、$r_{12}$ が $uv'-qvu'+kvv'=0$ を与える。
片側だけ $v$ が非零なら他側が最低点となり矛盾する。両 $v$ が零なら
$r_{02}$ から $Auu'=0$ となり、両側とも最低点である。
従ってこの点以外は $H=H'=1$ のアフィンチャートで尽きる。

このチャートでは $r_{23},r_{13},r_{12},r_{03}$ の順に $v',u',z',z$ を解く。
最初の二式はアフィン線型で、後二式から $z,z'$ は $u,v$ の二次以下の多項式となる
（$r_{12}$ の $z'$ 係数は $B\ne0$）。残る二式 $F,G$ は $u,v$ の三次式である。
従って共通成分がなければ有限で、アフィンにある全局所環の埋込み次元は高々2である。
最低点では $B\ne0$ により $r_{12}$ の一次部が $r_{01},r_{02},r_{03}$ の一次部から独立で、
残る二式の一次部は零なので接次元2である。

二平面三次式を斉次化し、無限遠で調べる。$[1:0:0]$ は全パラメータで横断交点である。
もう一方のチャート $[t:1:0]$ で $F,G$ の最高部はそれぞれ次数2と1であり、後者の単根を $r$ とする。
$-1$ 型では前者の $r$ での値が各辺で $1/2$ なのでこれ以外の無限遠交点はない。
従って共通曲線はなくアフィン長は $9-1=8$ である。
四乗根ではこの値が零になる。共通成分があれば無限遠で $r$ だけを通り、
$G$ の無限遠への制限のこの根が単純なので、その成分は直線 $u=rv+k$ でなければならない。
$G(rv+k,v)$ の二次係数が $k$ を一意に定める。その $k$ を代入した $F,G$ が
ともに恒等零となる条件の既約行階段消去（又はGröbner基底）は、全四辺で
ちょうど $(b,c,d,e,f)$ となる。前記コードは最高部・横断性・五式のイデアルを記号計算する。
これで例外の全てが決まる。例外上でも同じ階数6の再帰と良い還元から正則性が従う。
\end{proof}

\begin{corollary}\label{cor:finite-quantum-planar-lengths}
各正規形には長さが二十の半分を超える唯一の点がある。$-1$ 型での長さは12である。
原始四乗根では次表の行を上から読み、満たす最初の行がその長さを与える。
\[
\begin{array}{c|c}
\text{条件}&\text{長さ}\\\hline
 b\ne0&13\\
 d\ne0&14\\
 f\ne0&15\\
 ce\ne0&17\\
 e\ne0\text{ 又は }ac\ne0&18\\
 c\ne0&19
\end{array}
\]
\end{corollary}
\begin{proof}
前証明の第二無限遠点の交点数を $m$ とすればアフィン長は $8-m$、
元の窓の最低点の長さは $12+m$ である。斉次化の第三変数を $h$ とし、
$F(t,1,h),G(t,1,h)$ の $t$ に関する終結式の $h$ での位数が $m$ である。
$F(t,1,0)$ の二次の先頭係数は非零定数で $G(t,1,0)$ の零点は $r$ のみだからである。
前記コードでこの終結式の次数 $1,2,3,4,5$ までの逐次係数イデアルを取ると、
順に新しく $b,d,f,0,ce$ が加わる。次数6でさらに $ac^2-e^2,e^3$、
次数7でさらに $c^3$ が加わり、次数8で変わらない。有限点の条件の下で
これらの消失を順に読めば表になる。$-1$ 型はアフィン長8と全長20から直ちに従う。
\end{proof}
